%=======================================================================
\chapter{Optimisation des descripteurs (Phase~1)}

\section{Motivation}

L'efficacité d'un système de classification dépend fortement de la qualité des
caractéristiques utilisées pour représenter les données. Dans le cas des images IRM
cérébrales, les performances finales ne sont pas uniquement influencées par le choix du
classificateur, mais également par les paramètres utilisés lors de l'extraction des
caractéristiques.

Une approche classique consiste à évaluer directement chaque combinaison de paramètres à
travers l'entraînement complet d'un modèle de Machine Learning, généralement via une
procédure de type GridSearch associée à une validation croisée. Toutefois, cette stratégie
présente plusieurs limitations.

Premièrement, le nombre de combinaisons possibles devient rapidement très important lorsque
plusieurs descripteurs sont considérés simultanément. Deuxièmement, chaque évaluation
nécessite l'extraction complète des caractéristiques puis l'entraînement d'un ou plusieurs
modèles, ce qui entraîne un coût computationnel élevé. Enfin, les résultats obtenus peuvent
dépendre du classificateur utilisé, rendant difficile l'évaluation intrinsèque de la qualité
des caractéristiques.

Afin de réduire ce coût tout en conservant une évaluation pertinente, nous proposons une
phase d'optimisation indépendante des classificateurs. L'idée consiste à mesurer
directement la capacité discriminante des caractéristiques extraites à l'aide de métriques
statistiques de séparabilité. Cette approche permet d'identifier rapidement les
configurations les plus prometteuses avant d'effectuer le benchmark complet des modèles de
Machine Learning.

L'optimisation est réalisée séparément pour les descripteurs GLCM, LBP, DWT et HOG afin de
sélectionner les paramètres offrant la meilleure capacité de discrimination entre les
différentes classes de tumeurs cérébrales.

% -----------------------------------------------------------------------
\section{Métriques de séparabilité}

Pour évaluer la qualité des caractéristiques extraites, trois métriques complémentaires ont
été utilisées. Chacune mesure un aspect particulier de la séparation entre les classes dans
l'espace des caractéristiques. Les calculs sont effectués sur l'ensemble d'apprentissage
(\NTrain{} images) après extraction des descripteurs pour chaque configuration testée.

\subsection{Fisher Discriminant Ratio (FDR)}

Le Fisher Discriminant Ratio mesure la capacité d'une caractéristique à séparer différentes
classes en comparant la variance inter-classes à la variance intra-classe.

Une caractéristique est considérée comme discriminante lorsque les moyennes des différentes
classes sont éloignées tout en conservant une faible dispersion interne.

La formulation générale est donnée par~:
\begin{equation}
  \mathrm{FDR} = \frac{\sigma^2_{\mathrm{inter}}}{\sigma^2_{\mathrm{intra}}}
\end{equation}
où~:
\begin{itemize}[leftmargin=1.4em,itemsep=2pt]
  \item $\sigma^2_{\mathrm{inter}}$ représente la variance entre les classes~;
  \item $\sigma^2_{\mathrm{intra}}$ représente la variance à l'intérieur des classes.
\end{itemize}

Dans l'implémentation, le FDR est calculé pour chaque dimension du vecteur de
caractéristiques puis moyenné sur l'ensemble des descripteurs. Des valeurs élevées
indiquent une meilleure séparation entre les catégories de tumeurs.

\subsection{Mutual Information (MI)}

L'Information Mutuelle mesure la quantité d'information partagée entre une caractéristique et
les étiquettes de classes.

Contrairement au FDR, cette métrique ne suppose pas de relation linéaire entre les variables
et permet d'identifier des dépendances plus complexes. L'information mutuelle moyenne est
calculée entre chaque dimension et l'étiquette de classe à l'aide de
\texttt{mutual\_info\_classif}.

Une valeur élevée de Mutual Information indique qu'une caractéristique apporte une
information importante pour distinguer les différentes catégories de tumeurs.

\subsection{Davies--Bouldin Index (DBI)}

Le Davies--Bouldin Index est une mesure classique de qualité du clustering.

Il évalue simultanément~:
\begin{itemize}[leftmargin=1.4em,itemsep=2pt]
  \item la compacité des classes~;
  \item la séparation entre les groupes.
\end{itemize}

Contrairement aux deux métriques précédentes, un score faible indique une meilleure qualité
de séparation. Pour des raisons de coût computationnel en haute dimension, le DBI est
calculé après une réduction PCA à 20 composantes.

Dans notre procédure d'évaluation, le DBI est inversé après normalisation afin d'être
cohérent avec les autres métriques où une valeur élevée correspond à une meilleure
configuration.

% -----------------------------------------------------------------------
\section{Construction du score composite}

Chaque métrique fournit une vision partielle de la qualité des caractéristiques. Afin
d'obtenir une évaluation plus robuste, les trois indicateurs sont combinés dans un score
composite unique.

Pour chaque grille de recherche (GLCM, LBP, DWT ou HOG), les valeurs de FDR, MI et DBI sont
d'abord normalisées par \emph{Min--Max scaling} sur l'ensemble des configurations testées,
afin de les rendre comparables sur l'intervalle $[0,1]$. Le DBI normalisé est ensuite
inversé~:
\begin{equation}
  \mathrm{DBI}_{\mathrm{norm}} = 1 - \mathrm{MinMax}(\mathrm{DBI})
\end{equation}

Le score final est obtenu par combinaison pondérée à poids égaux~:
\begin{equation}
  \mathrm{Score}_{\mathrm{composite}}
  = w_1 \cdot \mathrm{FDR}_{\mathrm{norm}}
  + w_2 \cdot \mathrm{MI}_{\mathrm{norm}}
  + w_3 \cdot \mathrm{DBI}_{\mathrm{norm}}
\end{equation}
avec $w_1 = w_2 = w_3 = \tfrac{1}{3}$ dans notre implémentation.

Cette stratégie permet d'intégrer simultanément~:
\begin{itemize}[leftmargin=1.4em,itemsep=2pt]
  \item la séparation statistique des classes~;
  \item la pertinence informationnelle des caractéristiques~;
  \item la compacité des groupes.
\end{itemize}

Les configurations obtenant les scores composites les plus élevés sont retenues pour la phase
suivante du projet.

% -----------------------------------------------------------------------
\section{Recherche des meilleurs paramètres}

\subsection{Optimisation du descripteur GLCM}

Pour le descripteur GLCM, plusieurs configurations ont été explorées en faisant varier~:
\begin{itemize}[leftmargin=1.4em,itemsep=2pt]
  \item les distances entre pixels~;
  \item les orientations~;
  \item le nombre de niveaux de gris~;
  \item l'utilisation ou non de matrices symétriques.
\end{itemize}

Chaque configuration a été évaluée à l'aide du score composite précédemment défini.

Les résultats montrent que la meilleure configuration est obtenue avec~: \glcmParams{}
(score composite~: \glcmScore). Cette combinaison offre le meilleur compromis entre richesse
descriptive et capacité discriminante.

\begin{figure}[htbp]
  \centering
  \includegraphics[width=\textwidth]{glcm_param_search}
  \caption{Évolution du score composite pour les différentes configurations GLCM.}
  \label{fig:glcm-param-search}
\end{figure}

\subsection{Optimisation du descripteur LBP}

L'algorithme LBP a été optimisé en faisant varier~:
\begin{itemize}[leftmargin=1.4em,itemsep=2pt]
  \item le nombre de voisins ($P$)~;
  \item le rayon ($R$)~;
  \item la méthode de codage (\emph{uniform}, \emph{nri\_uniform}, \emph{ror}).
\end{itemize}

Après évaluation de l'ensemble des combinaisons testées, la meilleure configuration retenue
est~: \lbpParams{} (score composite~: \lbpScore). Cette configuration permet de capturer
efficacement les textures locales tout en limitant la sensibilité au bruit.

\begin{figure}[htbp]
  \centering
  \includegraphics[width=\textwidth]{lbp_param_search}
  \caption{Comparaison des scores composites obtenus pour les différentes configurations
    LBP.}
  \label{fig:lbp-param-search}
\end{figure}

\subsection{Optimisation du descripteur DWT}

Pour la Transformée en Ondelettes Discrètes, plusieurs familles d'ondelettes et niveaux de
décomposition ont été comparés parmi~: \texttt{haar}, \texttt{db1}, \texttt{db2},
\texttt{db4}, \texttt{sym2}, \texttt{coif1} et \texttt{bior1.3}.

Les expérimentations montrent que la configuration optimale correspond à~: \dwtParams{}
(score composite~: \dwtScore). Cette configuration fournit les caractéristiques
fréquentielles les plus discriminantes pour les images IRM du dataset étudié.

\begin{figure}[htbp]
  \centering
  \includegraphics[width=\textwidth]{dwt_param_search}
  \caption{Résultats de l'optimisation des paramètres DWT.}
  \label{fig:dwt-param-search}
\end{figure}

\subsection{Optimisation du descripteur HOG}

L'optimisation du descripteur HOG a porté sur~:
\begin{itemize}[leftmargin=1.4em,itemsep=2pt]
  \item le nombre d'orientations~;
  \item la taille des cellules (\emph{pixels per cell})~;
  \item la taille des blocs (\emph{cells per block})~;
  \item la méthode de normalisation (\emph{block norm}).
\end{itemize}

La meilleure configuration identifiée est~: \hogParams{} (score composite~: \hogScore).
Cette configuration permet d'obtenir une représentation robuste des contours et des
structures géométriques présentes dans les tumeurs cérébrales.

\begin{figure}[htbp]
  \centering
  \includegraphics[width=\textwidth]{hog_param_search}
  \caption{Influence des paramètres HOG sur le score composite.}
  \label{fig:hog-param-search}
\end{figure}

% -----------------------------------------------------------------------
\section{Synthèse de la Phase~1}

La phase d'optimisation a permis d'identifier les configurations offrant la meilleure
capacité de discrimination pour chaque famille de caractéristiques avant toute étape de
classification.

Le Tableau~\ref{tab:phase1-best} résume les paramètres optimaux retenus pour la suite du
projet.

\begin{table}[htbp]
  \centering
  \caption{Configurations optimales des descripteurs (Phase~1).}
  \label{tab:phase1-best}
  \begin{tabular}{@{}llc@{}}
    \toprule
    \thd{Descripteur} & \thd{Paramètres optimaux} & \thd{Score} \\
    \midrule
    GLCM & \glcmParams{} & \glcmScore \\
    LBP  & \lbpParams{}  & \lbpScore \\
    DWT  & \dwtParams{}  & \dwtScore \\
    HOG  & \hogParams{}  & \hogScore \\
    \bottomrule
  \end{tabular}
\end{table}

Cette première phase a permis de réduire l'espace de recherche tout en sélectionnant les
configurations les plus prometteuses. Les paramètres optimaux obtenus serviront de base à la
Phase~2, consacrée à la construction des ensembles de caractéristiques, à leur fusion et à
l'évaluation comparative des différents algorithmes de Machine Learning.
